\section{Red social de IA frente a redes sociales de internet}

A Elys le preguntaron si es una versión de Tinder con IA. La respuesta de Tristan es que no. El núcleo de Tinder es el emparejamiento de baja dimensión y la selección deslizando; lo que Elys quiere construir es un emparejamiento de context de alta dimensión y una socialización proactiva a través del doble digital. En las redes sociales tradicionales de internet, la persona tiene que expresarse, filtrar y romper el hielo por iniciativa propia; la socialización con IA busca que el doble digital se exprese y explore primero en tu lugar, y después te traiga de vuelta las posibles conexiones reales.

\begin{figure}[H]
\centering
\includegraphics[width=0.92\textwidth]{ai-social-vs-internet-social.png}
\caption{Diferencias entre las redes sociales de internet y la red social de IA.}
\end{figure}

\begin{knowledgebox}{Lectura de la figura: la diferencia de la socialización con IA no es solo «agregar IA»}
A la izquierda, las redes sociales de internet dependen de etiquetas de baja dimensión y del comportamiento proactivo del usuario; a la derecha, la socialización con IA depende de un context de alta dimensión y de la acción proactiva del doble digital. La diferencia real es que la IA reduce el costo de conectar, no que la interfaz se convierta en una ventana de chat.
\end{knowledgebox}

\subsection{Por qué no es un Tinder con IA}

Tinder se parece más a una criba rápida de candidatos para una relación, centrada en los datos visibles del perfil y en la elección inmediata. Elys quiere construir más bien una red de context: los valores, las aficiones y las formas de expresarse de una persona, por sutiles que sean o aunque no convenga decirlos en público, también pueden encontrar, mediante la IA, a alguien con quien resuenen. Lo que persigue es una conexión con menos fricción y de mayor dimensión.

\begin{warningbox}{Un error peligroso de la socialización con IA}
Si la socialización con IA solo genera una animación falsa, desgastará la confianza con rapidez. El usuario puede aceptar que la IA lo ayude a romper el hielo, pero el valor real sigue proviniendo de la conexión entre personas reales.
\end{warningbox}

\subsection{Quién posee más context}

En la entrevista se menciona que ChatGPT, Meta, Doubao (豆包) y otros podrían poseer una gran cantidad de context. El juicio de Tristan es que la inteligencia se irá distribuyendo de forma equitativa, pero el context no. Las empresas de modelos tienen la capacidad de los modelos, las plataformas sociales tienen las redes de relaciones, ChatGPT tiene la memoria personal y el historial de conversaciones, y las empresas emergentes de socialización con IA intentan construir productos especializados en torno al ciclo virtuoso del context.

\begin{figure}[H]
\centering
\includegraphics[width=0.92\textwidth]{context-competition.png}
\caption{Variables de competencia en la socialización con IA: capacidad de los modelos, redes de relaciones sociales, activos de Context y agencia del usuario.}
\end{figure}

\begin{knowledgebox}{Lectura de la figura: por qué el Context es más escaso que el modelo}
Los cuatro tipos de actores de la figura tienen cada uno su ventaja: las empresas de modelos tienen la inteligencia, las plataformas sociales tienen las relaciones, las empresas emergentes tienen un nuevo paradigma y el propio usuario tiene la agencia. Si la inteligencia se convierte poco a poco en una mercancía, el context de largo plazo y las relaciones entre personas reales pasarán a ser los activos más escasos.
\end{knowledgebox}

\subsection{Resumen de la sección}

El núcleo de la red social de IA no es una nueva generación de algoritmos de emparejamiento, sino el paso de una socialización basada en etiquetas de baja dimensión a una socialización basada en context de alta dimensión. Su objetivo es que la IA reduzca el costo de conectar a personas reales, no sustituir esa conexión.
